% Compile twice with: pdflatex Direct_interpretation_of_R_in_PTS_rk.tex
\documentclass[11pt,a4paper]{article}
\usepackage[T1]{fontenc}
\usepackage{lmodern,amsmath,amssymb,amsthm,mathtools,booktabs,microtype}
\usepackage[margin=25mm,headheight=14pt]{geometry}
\usepackage{fancyhdr,needspace}
\usepackage[hidelinks]{hyperref}
\hypersetup{pdftitle={A direct interpretation of Tarski--Mostowski--Robinson arithmetic R in PTS-rk},pdfsubject={A fixed ground-faithful arithmetic interpretation using certified capped arithmetic algebras}}
\pagestyle{fancy}\fancyhf{}
\fancyhead[L]{\small A direct interpretation of $\mathsf R$ in $\mathsf{PTS}_{rk}$}
\fancyhead[R]{\small\thepage}
\setlength{\parindent}{1.2em}\setlength{\parskip}{2pt}
\emergencystretch=1.5em
\newcommand{\rk}[1]{\mathsf{rk}_{#1}}
\newcommand{\PTS}{\mathsf{PTS}_{rk}}
\newcommand{\RR}{\mathsf R}
\newcommand{\Good}{\mathsf{Good}}
\newcommand{\Case}{\mathsf{Case}}
\newcommand{\Inv}{\mathsf{Inv}}
\newcommand{\zz}{\mathsf z}
\newcommand{\ssu}{\mathsf s}
\newcommand{\addD}{\mathsf a}
\newcommand{\mm}{\mathsf m}
\newcommand{\num}[1]{\nu_{#1}}
\newcommand{\ev}{\operatorname{ev}}
\newcommand{\nc}{\mathop{\neg}\nolimits_c}
\newcommand{\ac}{\mathbin{\wedge_c}}
\newcommand{\oc}{\mathbin{\vee_c}}
\newcommand{\Orc}{\mathop{\bigvee_{\!c}}\displaylimits}
\newcommand{\allA}{\Pi\mathcal A.\,}
\newtheorem{theorem}{Theorem}
\newtheorem{lemma}[theorem]{Lemma}
\newtheorem{proposition}[theorem]{Proposition}
\theoremstyle{remark}\newtheorem*{remark}{Remark}
\title{\bfseries A direct interpretation of\newline Tarski--Mostowski--Robinson arithmetic $\mathsf R$\newline in $\mathsf{PTS}_{rk}$}
\author{}\date{12 October 2026}
\begin{document}
\begin{center}
{\Large\bfseries A direct interpretation of\\[3pt]
Tarski--Mostowski--Robinson arithmetic $\mathsf R$\\[3pt]
in $\mathsf{PTS}_{rk}$\par}
\vspace{9pt}
12 October 2026
\end{center}
\thispagestyle{plain}
\begin{abstract}
We construct a fixed, ground-faithful interpretation of classical
Tarski--Mostowski--Robinson arithmetic $\RR$ in the rank-indexed pure type
system $\PTS$. Arithmetic expressions are Church encoded and equipped with
a bounded-rank induction certificate. Certified arithmetic observations
interpret equality and order. A single bounded case lemma verifies the two
quantified axiom schemes of $\RR$; finite capped arithmetic algebras separate
standard numerals. The construction uses only dependent products and
$\beta$-conversion, without an intermediate set or multiset theory.
\end{abstract}

\section{The target theory and the proof types}
The pure type system $\PTS$ has sorts $\rk i$ for external $i\in\mathbb N$,
sort axioms $\rk i:\rk0$, and the product rule
\[
 \frac{\Gamma\vdash A:\rk a\qquad\Gamma,x:A\vdash B:\rk b}
 {\Gamma\vdash\Pi x:A.B:\rk{\max(a+1,b)}}.
\]
The other rules are the ordinary PTS rules, with $\beta$-conversion only.
There is no cumulativity. Arrows associate to the right.

We use $\RR$ in the signature $\{0,S,+,\cdot,\leq\}$ with
$\bar n=S^n0$ and the following schemes, for external $m,n\in\mathbb N$
\cite{tmr}\cite[Definition~1.1]{kurahashi}:
\begin{align}
 \bar m+\bar n&=\overline{m+n}, \tag{R1}\label{eq:r1}\\
 \bar m\cdot\bar n&=\overline{mn}, \tag{R2}\label{eq:r2}\\
 \bar m&\ne\bar n\quad(m\ne n), \tag{R3}\label{eq:r3}\\
 \forall x\bigl(x\leq\bar n&\to\bigvee_{i=0}^n x=\bar i\bigr),
 \tag{R4}\label{eq:r4}\\
 \forall x\bigl(x\leq\bar n&\lor\bar n\leq x\bigr).
 \tag{R5}\label{eq:r5}
\end{align}
Source quantifiers will be relativized to a fixed predicate $\Good$.
Equality is a definable equivalence relation, not $\beta$-conversion.

Put
\[
 \bot:=\Pi Z:\rk0.Z:\rk1,\qquad
 \top:=\rk0\to\rk0:\rk1,\qquad \mathsf{id}:\top.
\]
Section~\ref{sec:empty} justifies that $\bot$ has no closed inhabitant.
Use the fixed-answer negative connectives
\begin{gather*}
 \nc A:=A\to\bot,\qquad
 A\ac B:=(A\to B\to\bot)\to\bot,\\
 \Orc_{i=0}^k A_i:=(A_0\to\bot)\to\cdots\to(A_k\to\bot)\to\bot.
\end{gather*}
A type $A$ is \emph{stable} if $((A\to\bot)\to\bot)\to A$ is inhabited.
Negative types are stable, as are dependent products with stable codomain.
The displayed conjunction and disjunction are stable and have the classical
rules when their component proof types are stable. They eliminate into
stable targets; in particular, case branches may all return $\bot$.
Implication is an arrow. No elimination of $\bot$ into arbitrary positive
types is assumed.

For exact-rank motives define $J_1:=\top$, $J_{k+1}:=J_k\to J_k$, with
canonical inhabitants $j_k$. If $A:\rk a$, $a\leq k$, and $k\geq2$, then
\[
 L_k(A):=J_{k-1}\to A:\rk k.
\]
Abstraction and evaluation at $j_{k-1}$ give maps in both directions.
We use this lift whenever a motive has rank below the specified rank.

\section{Arithmetic expressions and certified observations}
\subsection*{The fixed domain}
Use the Church presentation type
\[
 D:=\Pi X:\rk1.\ X\to(X\to X)\to(X\to X\to X)
                         \to(X\to X\to X)\to X:\rk3.
\]
Its four constructors are the following closed lambda terms:
\begin{align*}
 \zz&:=\lambda X\,z\,s\,a\,m.\ z,\\
 \ssu d&:=\lambda X\,z\,s\,a\,m.\ s(d\,X\,z\,s\,a\,m),\\
 \addD d e&:=\lambda X\,z\,s\,a\,m.\ a(d\,X\,z\,s\,a\,m)(e\,X\,z\,s\,a\,m),\\
 \mm d e&:=\lambda X\,z\,s\,a\,m.\ m(d\,X\,z\,s\,a\,m)(e\,X\,z\,s\,a\,m).
\end{align*}
Binder annotations are dictated by $D$. Set $\num n:=\ssu^n\zz$.
These are \emph{arithmetic-expression presentations}: $\addD$ and $\mm$ build
syntax. Their ground arithmetic laws will hold in interpreted equality;
they are not asserted to reduce to numerals by $\beta$ alone.

Define the domain predicate by
\begin{align*}
 \Good(d):={}&\Pi Q:D\to\rk{10}.\ Q(\zz)\\[-2pt]
 &\to(\Pi u:D.\ Q(u)\to Q(\ssu u))\\[-2pt]
 &\to(\Pi u,v:D.\ Q(u)\to Q(v)\to Q(\addD u v))\\[-2pt]
 &\to(\Pi u,v:D.\ Q(u)\to Q(v)\to Q(\mm u v))\to Q(d).
\end{align*}
This has rank $12$. It is a definition, not an induction axiom on all of $D$.
The four constructors preserve $\Good$ by direct abstraction. For example,
from $g:\Good(d)$ the successor certificate is
\[
 \lambda Q\,q_0\,q_s\,q_a\,q_m.\
 q_s\,d\,(g\,Q\,q_0\,q_s\,q_a\,q_m).
\]
The binary certificates apply $q_a$ or $q_m$ to the two child certificates.
Hence the domain is nonempty and closed under every interpreted function.

\subsection*{The observation telescope}
An observation has parameters
\[
 X:\rk1,\quad z:X,\quad s:X\to X,\quad a,m:X\to X\to X,
 \quad\delta:X\to\rk4,\quad R_0,C_0:X\to X\to\rk1.
\]
Abbreviate the \emph{negative} proof types
\[
 u\simeq v:=R_0(u,v)\to\bot,\qquad
 u\preceq v:=C_0(u,v)\to\bot.
\]
Both have rank $2$. The telescope also contains separate proofs of the
following laws. Every object variable is guarded by $\delta$: thus a law
$H(u,v)$ means $\Pi u:X.\delta(u)\to\Pi v:X.\delta(v)\to H(u,v)$.
All displayed paired laws mean separate proof arguments, not encoded records.

\emph{Domain and congruence.} The domain contains $z$ and is closed under
$s,a,m$. On it, $\simeq$ is an equivalence relation respected by $s,a,m$;
$\preceq$ respects $\simeq$ in both arguments and is reflexive and transitive.
For example, its congruence certificate concludes
\[
 u\simeq u'\to v\simeq v'\to u\preceq v\to u'\preceq v'.
\]
\emph{Arithmetic equations.}
\begin{align}
 a(z,u)&\simeq u,&a(su,v)&\simeq s(a(u,v)),\notag\\
 m(z,u)&\simeq z,&m(su,v)&\simeq a(v,m(u,v)),&m(u,z)&\simeq z.
 \label{eq:alg-equations}
\end{align}
\emph{Order propagation}, with $1_X:=sz$:
\begin{align}
 u&\preceq su,&u&\preceq a(u,v),&v&\preceq a(u,v),\notag\\
 1_X\preceq v&\to u\preceq m(u,v),&&
 1_X\preceq u\to v\preceq m(u,v).
 \label{eq:alg-order}
\end{align}
No successor injectivity, cancellation, or unbounded induction is required.
Each individual certificate has rank at most $5$: after its domain guards
are opened, a relation has rank $2$ and an implication between relations
has rank at most $3$.

Write $\Pi\mathcal A$ for this fixed parameter telescope followed by these
individual proof arguments. Define
\[
 \ev_{\mathcal A}(d):=d\,X\,z\,s\,a\,m
\]
and
\begin{align*}
 E(d,e)&:=\allA\ev_{\mathcal A}(d)\simeq\ev_{\mathcal A}(e),\\
 O(d,e)&:=\allA\ev_{\mathcal A}(d)\preceq\ev_{\mathcal A}(e),\\
 G(d)&:=\allA\delta(\ev_{\mathcal A}(d)).
\end{align*}
All three have rank $6$; $E,O$ are stable, while $G$ is positive.
The certificates are quantified assumptions of an observation, \emph{not}
axioms added to $\PTS$. The next section constructs closed instances.

\begin{lemma}\label{lem:basic}
$G$ contains $\zz$ and is closed under $\ssu,\addD,\mm$, and
$\Pi d:D.\Good(d)\to G(d)$ is inhabited. On $G$-elements, $E$ is an
equivalence, all three operations respect $E$, and $O$ is a congruent
preorder. Equations~\eqref{eq:alg-equations} and
\eqref{eq:alg-order} hold globally with $E,O$.
\end{lemma}
\begin{proof}
The fold equations are beta equations. They and the domain certificates
prove closure of $G$. Apply $\Good(d)$ to $u\mapsto L_{10}(G(u))$ and
then evaluate at $j_9$. Every remaining claim follows by opening the
observation telescope and supplying the relevant $G$-proofs as guards.
\end{proof}
External induction on the natural numbers in a \emph{closed} arithmetic
calculation now constructs proofs of
\[
 E(\addD\num m\num n,\num{m+n}),\quad
 E(\mm\num m\num n,\num{mn}),\quad
 O(\num i,\num j)\quad(i\leq j).
\tag{*}\label{eq:ground}
\]
For the first two use the recursive equations; for the last use reflexivity,
$u\preceq su$, and transitivity. No internal arithmetic induction is used.

\section{Finite arithmetic observations and separation}
For each external $K\geq0$, put $c=K+1$ and take the finite chain
$C_K=\{0,\ldots,c\}$, with
\[
 z=0,\quad s(i)=\min(i+1,c),\quad
 a(i,j)=\min(i+j,c),\quad m(i,j)=\min(ij,c).
\]
Equality and order are the actual equality and order of this chain. All the
observation laws hold. For the multiplication recursion, $j=0$ gives zero
on both sides; if $j>0$, reaching the cap makes both sides equal $c$.
For multiplication propagation, $1\leq j$ gives
$\min(ij,c)\geq i$, and symmetrically. Addition propagation and the remaining
laws follow directly from the chain operations. In particular, capped
multiplication still satisfies $c\cdot0=0$.

\begin{lemma}\label{lem:finite}
Every $C_K$ has a closed implementation of the entire observation telescope.
Consequently $E(\num m,\num n)\to\bot$ is inhabited when $m\ne n$, and
$O(\num m,\num n)\to\bot$ is inhabited when $m>n$.
\end{lemma}
\begin{proof}
Put $q=K+2$ and use the fixed-rank selector type
\[
 F_q:=\underbrace{\rk1\to\cdots\to\rk1}_{q\text{ arguments}}\to\rk1:\rk1,
 \qquad e_i:=\lambda Z_0\ldots Z_{q-1}.Z_i.
\]
A unary table $h$ and a binary table $b$ have the realizations
\begin{align*}
 \widehat h(u)&:=\lambda\vec Z.\ u\,Z_{h(0)}\cdots Z_{h(q-1)},\\
 \widehat b(u,v)&:=\lambda\vec Z.\
 u\,(v\,Z_{b(0,0)}\cdots Z_{b(0,q-1)})\cdots
       (v\,Z_{b(q-1,0)}\cdots Z_{b(q-1,q-1)}).
\end{align*}
Interpret $z$ by $e_0$ and the operations by their tables. For a relation
with truth table $T_{ij}$, replace $Z_{b(i,j)}$ in the binary formula by
$B_{ij}$, where $B_{ij}=\bot$ for a true entry and $B_{ij}=\top$ for a false
one, and omit the outer $\lambda\vec Z$. This defines $R_0,C_0$ with values
of rank $1$. Their negations have exactly the required truth values:
$\bot\to\bot$ has its identity proof, whereas $\top\to\bot$ yields
$\bot$ when applied to $\mathsf{id}:\top$.

Use the domain certificate
\[
 \delta_q(u):=\Pi Q:F_q\to\rk3.\
 Q(e_0)\to\cdots\to Q(e_{q-1})\to Q(u):\rk4.
\]
Each selector has its projection certificate. For operation closure, open
the output $\delta_q$ and eliminate the input certificates into its
rank-$3$ motive $Q$ composed with the operation. For relation laws, open
all domain guards and use selector elimination into the remaining goal of
rank at most $3$, lifted to rank $3$ if necessary. This reduces the proof
to the finite table checks above. A false relational premise gives $\bot$;
every relational conclusion is negative, so its proof follows by abstraction.
There is no assertion that every raw inhabitant of $F_q$ is a selector.

For separation choose $K\geq\max(m,n)$. Both folds are their distinct
selectors, or occur in strictly reversed order. Specializing a putative
$E$- or $O$-proof to this closed observation gives $\top\to\bot$; apply it
to $\mathsf{id}$. For example, the same argument at cutoff $K=n$ refutes
$O(\num{n+1},\num n)$.
\end{proof}

\section{One bounded numerical case lemma}
For each fixed external $n$, define the flat negative disjunction
\[
 \Case_n(d):=\Orc\bigl(E(d,\num0),\ldots,E(d,\num n),O(\num{n+1},d)\bigr).
\]
It has rank $8$, independent of $n$. The final branch is a \emph{global
lower bound}, not just an observation that a finite classifier overflowed.
This eliminates the need for a reconstruction theorem.

\begin{lemma}\label{lem:cases}
For every external $n\geq0$, the type
$\Pi d:D.\Good(d)\to\Case_n(d)$ is inhabited.
\end{lemma}
\begin{proof}
Retain the positive domain evidence using the rank-$10$ motive
\[
 \Inv_n(d):=\Pi Z:\rk8.\bigl(G(d)\to\Case_n(d)\to Z\bigr)\to Z.
\]
We construct its four closure proofs. In each constructor step, open the
parent's $Z$ and continuation and eliminate the child invariants at that
\emph{same} $Z$. They provide the positive $G$-proofs and the negative case
proofs separately. Lemma~\ref{lem:basic} constructs the parent's $G$-proof
without any case elimination. To construct its case proof, open its branch
continuations; all remaining branches have answer type $\bot$. The case
proofs can now be split. The following deductions use only
Lemma~\ref{lem:basic} and the closed facts~\eqref{eq:ground}.

\emph{Zero and successor.} Zero has its reflexive $E$-proof. A successor of
an $i$-case is the $(i+1)$-case if $i<n$, and supplies
$O(\num{n+1},\ssu d)$ if $i=n$. From the last branch use
$O(\num{n+1},d)$ and $O(d,\ssu d)$.

\emph{Addition.} Two finite cases $i,j$ give equality to $\num{i+j}$.
If $i+j\leq n$, take that finite branch; otherwise the ground comparison
$n+1\leq i+j$ gives the final branch. If either child has the final branch,
use the corresponding bound of that child by the sum.

\emph{Multiplication.} Two finite cases give equality to $\num{ij}$ and are
handled as above. If either factor is equal to zero, use the appropriate
zero equation, including $m(u,z)\simeq z$. Otherwise a finite factor has
index at least one. A factor in the final branch is also at least one,
because $1\leq n+1$. Thus, if one factor has the final branch and the other
is not in the zero case, multiplication propagation bounds the product
below by the large factor. This supplies its final branch. These alternatives
cover both finite/large cases and the large/large case.

Apply the opened parent continuation to its positive $G$-proof and the
constructed case proof. This proves closure of $\Inv_n$ without eliminating
negative cases into a positive type. Now instantiate $\Good(d)$ directly
at $\Inv_n:D\to\rk{10}$. Project $\Case_n(d)$ by choosing
$Z=\Case_n(d)$ and continuation $\lambda g\,c.c$.
\end{proof}

The ranks used in the construction are
\[
\begin{array}{c|ccccccc}
\text{type}&D&\delta_q&\text{algebra certificates}&E,O,G&\Case_n&\Inv_n&\Good\\
\hline
\text{rank}&3&4&\leq5&6&8&10&12.
\end{array}
\]
The bound is independent of the numerical values and table sizes; no
minimality of these ranks is claimed.

\section{Verification of the interpretation}
Interpret $0,S,+,\cdot$ by $\zz,\ssu,\addD,\mm$, and interpret the atoms
$t=u$ and $t\leq u$ by $E(t^\dagger,u^\dagger)$ and
$O(t^\dagger,u^\dagger)$. Use the negative connectives of Section~1 and
\[
 (\forall x\,\varphi)^\dagger
 :=\Pi x:D.\Good(x)\to\varphi^\dagger(x),\qquad
 (\exists x\,\varphi)^\dagger
 :=\bigl(\Pi x:D.\Good(x)\to\varphi^\dagger(x)\to\bot\bigr)\to\bot.
\]
All translated formulas are stable. Thus the classical first-order rules
are admissible. Domain closure and Lemma~\ref{lem:basic} supply the equality
and function obligations. Invariance of $\Good$ on arbitrary raw
$E$-equivalent presentations is not required: equality substitution is
proved for admitted arguments.

\begin{theorem}\label{thm:main}
These fixed definitions interpret classical $\RR$ in $\PTS$. Every
translated axiom has an effectively constructible closed proof term.
The interpretation is faithful on ground arithmetic equations and comparisons:
true instances are inhabited and false instances have internal refutations.
\end{theorem}
\begin{proof}
Schemes (R1) and (R2) are~\eqref{eq:ground}; (R3) is
Lemma~\ref{lem:finite}.

For (R4), fix external $n$, and assume $d:D$, $g:\Good(d)$ and
$h:O(d,\num n)$. Apply Lemma~\ref{lem:cases}. A finite branch already gives
one of $E(d,\num0),\ldots,E(d,\num n)$. The final branch together with $h$
gives $O(\num{n+1},\num n)$ by transitivity, contradicting
Lemma~\ref{lem:finite}. The required finite disjunction is stable, so this
case analysis is valid.

For (R5), use the same cases. In an $i$-branch, the ground fact $i\leq n$
and congruence give $O(d,\num n)$. In the final branch, compose
$O(\num n,\num{n+1})$ with $O(\num{n+1},d)$ to obtain $O(\num n,d)$.
Introduce the appropriate side of the translated disjunction.

Finally, external evaluation of any closed arithmetic expression reduces
it, in $E$, to its numeral by~\eqref{eq:ground} and congruence. Apply the
same results and Lemma~\ref{lem:finite} to its equality and order comparisons.
All external inductions and finite tables produce finite lambda terms;
the domain, operations, relations, and translation are fixed independently of
the axiom instance. No conservativity claim for arbitrary first-order
sentences is made.
\end{proof}

\section{Nontriviality of the answer type}\label{sec:empty}
Only the uninhabitation of $\bot$ is needed to ensure that the interpretation
is nonvacuous. Here is a short relative set-model justification, in external
$\mathsf{ZFC}+\mathsf{GCH}$; this metatheory adds no axiom to $\PTS$.
Choose distinct codes $c_\kappa$ for all cardinals below some finite-index
aleph, and distinct codes $s_j$ for $j\in\mathbb N$. Put
\[
 U_i=\{c_\kappa:\kappa\leq\aleph_i\}\cup\{s_j:j\in\mathbb N\},\qquad
 \operatorname{El}(c_\kappa)=\kappa,\quad
 \operatorname{El}(s_j)=U_j.
\]
Interpret $\rk i$ by $s_i$. Every $U_i$ is countable and $s_i\in U_0$.
For $a\in U_m$ and $b(x)\in U_n$, the full dependent function set has size
\[
 \left|\prod_{x\in\operatorname{El}(a)}\operatorname{El}(b(x))\right|
 \leq\aleph_n^{\aleph_m}\leq\aleph_{\max(m+1,n)}.
\]
Choose its cardinal code and a decoding bijection, once for each actual pair
$(a,b)$, independently of the typing derivation or rank bound. All eligible
pairs form a set. Application uses the chosen bijection and abstraction its
inverse; pointwise reindexing gives substitution and beta soundness. This
interprets exactly the stated product and sort rules. The product decoding
$\Pi Z:\rk0.Z$ has an empty factor at $c_0\in U_0$, so is empty. Consequently
there is no closed proof of $\bot$, as used above.

\begin{thebibliography}{2}\small
\bibitem{tmr}
A. Tarski, A. Mostowski, and R. M. Robinson.
\emph{Undecidable Theories}. North-Holland, 1953.
\bibitem{kurahashi}
T. Kurahashi.
\emph{Incompleteness and undecidability of theories consistent with
$\mathsf R$}. arXiv:2211.15455v1, 2022, Definition~1.1.
\url{https://arxiv.org/abs/2211.15455v1}.
\end{thebibliography}
\end{document}
